\documentclass[a4paper]{article}
\usepackage{amsfonts}
\usepackage{amsmath}
\usepackage{amssymb}
\usepackage{graphicx}

\begin{document}

\noindent \large Auteur: Rob Willemsen \\
\noindent \large Studierichting: Informatica\\
\noindent \large Oefeningenreeks 1C nr. 31.11\\

\medskip
\normalsize

Gegeven is de macht $z = (-\sqrt{3} - i)^{20}$. \\

\textbf{Basisgetal omzetten naar poolvorm:} \\

Stel $w = -\sqrt{3} - i$. \\

$r = \sqrt{(-\sqrt{3})^2 + (-1)^2} = \sqrt{3 + 1} = 2$ \\

$\tan \alpha = \dfrac{-1}{-\sqrt{3}} = \dfrac{\sqrt{3}}{3}$ \\

$\alpha = \pi + \dfrac{\pi}{6} = \dfrac{7\pi}{6}$ \\

$w = 2 \operatorname{cis} \dfrac{7\pi}{6}$ \\
 
\textbf{Macht verheffen met De Moivre:} \\

$z = w^{20} = \left(2 \operatorname{cis} \dfrac{7\pi}{6}\right)^{20}$ \\

$z = 2^{20} \operatorname{cis} \left(20 \cdot \dfrac{7\pi}{6}\right)$ \\

$z = 2^{20} \operatorname{cis} \left(\dfrac{140\pi}{6}\right) = 2^{20} \operatorname{cis} \left(\dfrac{70\pi}{3}\right)$ \\

\textbf{Argument vereenvoudigen en uitwerken:} \\

$\dfrac{70\pi}{3} - \dfrac{66\pi}{3} = \dfrac{4\pi}{3}$ \\

$z = 2^{20} \operatorname{cis} \dfrac{4\pi}{3}$ \\

$z = 2^{20} \left(\cos \dfrac{4\pi}{3} + i \sin \dfrac{4\pi}{3}\right)$ \\

$z = 2^{20} \left(-\dfrac{1}{2} - i\dfrac{\sqrt{3}}{2}\right)$ \\

$z = -2^{19} - 2^{19}i\sqrt{3}$ \\

\textbf{Antwoord:} $z = -2^{19} - 2^{19}i\sqrt{3}$ \\

\begin{thebibliography}{999}
%\bibitem{a} Referentie
\end{thebibliography}
\end{document}